This paper follows a strict rule throughout: every threshold, cutoff, or
definition used in the analysis was chosen and written down \emph{before}
its effect on any result was seen, together with the reasoning behind it
and a plan for testing how sensitive the results are to it. Where a choice
was later revised, the earlier version and the reason for the change are
both kept on record rather than quietly edited away. This section explains
every such choice in plain terms.

\subsection{Study period and phases}
\label{sec:phases}

The study covers \PZeroStart{} through \PFourBoundary{}. For analysis, this
period is split into five phases, summarized in Table~\ref{tab:phases}.
The boundaries between phases were not taken from press reports or public
timelines; they were found directly in the data, by scanning the entire
study period for sustained drops in Iran's country-level network
responsiveness (at least two days long, using the same responsiveness
threshold defined in Section~\ref{sec:state-derivation}) without assuming in
advance where a shutdown might be. This scan surfaced exactly three such
drops: a short one in June 2025 (a previously known reference event, used
here as a smaller-scale comparison case), an 18-day one beginning
\POneStart{} that had not previously been identified as a distinct event,
and the long shutdown from late February through May 2026. Notably, the
scan found no drop at all around the commonly reported ``late December 2025''
start date for a gradual throttling period -- every signal checked stays
flat and normal through all of December 2025, so this paper's baseline
period simply extends up to the first point where a real change is
detectable in the data, rather than stopping at a press-reported date that
the measurement does not support.

\begin{table}[t]
\centering
\small
\caption{The five study phases. All boundaries are UTC.}
\label{tab:phases}
\begin{tabular}{@{}p{0.6cm} p{1.6cm} p{4.2cm} p{6.5cm}@{}}
\toprule
Phase & Name & Period & What happened \\
\midrule
P0 & \PZeroName{} & \PZeroStart{} -- \PZeroBoundary{} & Normal conditions, including a brief reference blackout in June 2025. \\[4pt]
P1 & \POneName{} & \POneStart{} -- \POneBoundary{} & An 18-day near-total blackout, followed by roughly a month back to normal before the main shutdown. \\[4pt]
P2 & \PTwoName{} & \PTwoStart{} -- \PTwoBoundary{} & The main shutdown begins; connectivity collapses within about 40 minutes. \\[4pt]
P3 & \PThreeName{} & \PThreeStart{} -- \PThreeBoundary{} & Sustained, near-total blackout. \\[4pt]
P4 & \PFourName{} & \PFourStart{} -- \PFourBoundary{} & Partial, uneven recovery. \\
\bottomrule
\end{tabular}
\end{table}

The Prelude phase (P1) is a single labeled span but is not uniform: it
contains the 18-day blackout mentioned above followed by a quieter
month-long stretch before the main shutdown begins. Results for this phase
are always reported with the 18-day event called out on its own, rather
than averaged together with the quieter stretch that follows it.

The exact instant the main shutdown begins (the start of the Onset phase,
P2) is measured to the minute using the network-responsiveness signal,
which is recorded roughly every ten minutes -- far finer than the 8-hour
snapshots used for the routing measurements described below. This matters
because the Onset phase itself is short (about 16 hours) and only two
routing snapshots fall inside it, so any result reported for this phase
individually is really an average of just two measurements. Checking that
average against nearby snapshots shows it is a stable, representative
number, not a fragile average of only two points: connectivity drops
sharply between \trajTimePrevMidnight{} and \trajTimeTransitionMorning{}
(from \trajDarkSharePrevMidnight{} of networks affected to
\trajDarkShareTransitionMorning{}), and stays at essentially that same level
through \trajTimeTransitionAfternoon{} (\trajDarkShareTransitionAfternoon{})
and into the following phase (\trajDarkShareNextMidnight{}) -- a single
sharp step that immediately settles at the level the shutdown holds for
months afterward, not a number that would have looked very different had
the boundary been drawn an hour or two earlier or later.

\subsection{Network population and classification}

The set of Iranian networks studied (\nIrAsns{} networks, \nIrPrefixes{}
address blocks) is derived from public Internet registry records, which
list which address blocks are assigned to which country and organization.

Each network is assigned one of ten categories -- state telecom, mobile,
internet service provider (ISP), government, financial, energy, media,
hosting, education, or other -- based on the service it actually provides,
not who owns it. A state-owned consumer internet provider is categorized as
an ordinary ISP, for example; the state-telecom category is reserved for
the core operators that carry Iran's international connectivity,
regardless of their corporate structure. Category assignment combines
three approaches: the largest 100 or so networks by address space, every
network that shows up in a network-behavior group relevant to the
selectivity question this paper asks (Section~\ref{sec:h3}), and a
name-matching pass over official registry organization names (for example,
matching ``bank,'' ``finance,'' or ``payment'' to the financial category) that
catches small but analytically important sectors -- banks and payment
processors typically use much less address space than large telecom
operators, so a pure size cutoff would have missed exactly the sectors this
paper's selectivity question is about. Name matches only suggest
candidates; every category assignment is a human judgment call, checked
against public information about each organization. Categories with fewer
than 20 classified networks (mobile: \nTypeMobile{}; education:
\nTypeEducation{}) are reported only as raw counts, never percentages,
since percentages of that few networks are not meaningful. \nClassified{}
of \nIrAsns{} networks have been categorized this way; the rest are left
uncategorized and appear only in the ``unclassified'' group in
category-based results.

Category assignments were proposed by researching each organization's
publicly available information, then independently checked: a second
reviewer, without seeing the proposed categories, blind-coded a random
sample of \kappaN{} networks. The two sets of categorizations agreed on
\kappaAgreePct{} of cases (Cohen's $\kappa$ = \kappaValue{}, comfortably
above the 0.7 threshold set in advance for judging the categorization
reliable).

\subsection{Control population}

To distinguish Iran-specific events from artifacts of the measurement
setup itself, the same signals are also tracked for \nControlAsnsTotal{}
networks in three neighboring countries with no documented nationwide
shutdown during the study period (Turkey, the United Arab Emirates, and
Pakistan; \nControlAsnsPerCountry{} networks each). These control networks
were selected mechanically -- the largest networks by address space in
each country that meet the same data-adequacy bar used for Iran (below) --
and fixed in advance, before being used to check any result. A given time
window is flagged as a possible measurement artifact if the same kind of
anomaly (a responsiveness drop or a routing-visibility drop) shows up in at
least \controlArtifactBinShare{} of the control networks at the same time.
Every anomaly reported for Iran in this paper is checked against this
control group, and anything that shows up equally in the controls is
treated as a measurement artifact, not a real Iranian event.

\subsection{Measuring routing visibility}
\label{sec:visibility}

Global Internet routers continuously exchange messages announcing which
address blocks they can reach, via the Border Gateway Protocol (BGP). A
subset of these routers, called route collectors, publicly and
continuously record everything they see; two long-running collector
projects, RouteViews and the RIPE Routing Information Service (RIS), are
the data source for this paper's routing measurements. If a network's
address block is being announced and reaching these collectors, the block
is ``visible''; if collectors stop hearing it announced at all, it has been
``withdrawn.''

Visibility is measured every \ribIntervalHours{} hours across the full
study period, using a fixed set of collectors (RouteViews' two most
complete vantage points) as the main data source. A collector counts as
providing full coverage at a given time only if it is carrying at least
\fullFeedMinIpvFour{} IPv4 routes or \fullFeedMinIpvSix{} IPv6 routes; the
visibility measurement at each point in time only counts collectors
meeting this bar at that specific time, rather than assuming a fixed set of
collectors are always reliable. This matters because collector data feeds
occasionally have real gaps or outages that have nothing to do with events
in Iran, and without this check such a gap could easily be mistaken for an
Iranian connectivity event. Time points where too few collectors meet this
bar (fewer than \minFullfeedPeers{} -- a level chosen so that the
statistical noise in the measurement stays small relative to the
thresholds used elsewhere in this paper) are excluded from
visibility-based results rather than reported on unreliable data.

A second, supplementary set of measurements uses a wider range of
collectors (RouteViews and RIS combined, several hundred vantage points)
but only during the periods this paper's questions are most sensitive to
vantage-point choice: the main shutdown's onset, the entire recovery phase,
one baseline reference month, and three sample weeks during the sustained
blackout. This wider measurement is used only for the specific secondary
checks noted where it appears, never mixed with the main measurement in a
single comparison, and exists to confirm the main results are not an
artifact of which collectors happen to be used. That confirmation is
direct, not assumed: comparing the two measurements at the same points in
time, \bimodalNearOnePrimary{} of values in the main measurement are close
to fully visible and \bimodalNearZeroPrimary{} are close to fully
withdrawn, with only \bimodalAmbiguousPrimary{} landing anywhere in
between; the wider measurement shows an almost identical split
(\bimodalNearOneRis{} / \bimodalNearZeroRis{} / \bimodalAmbiguousRis{})
despite drawing on a substantially larger and different set of vantage
points. In other words, visibility in this data is close to a clean
on/off signal, and that stays true regardless of which collector set is
used -- which is why a single fixed threshold for ``visible'' works well
(Section~\ref{sec:state-derivation}), as Figure~\ref{fig:bimodality} shows
directly.

\begin{figure}[t]
  \centering
  \includegraphics[width=0.85\linewidth]{figures/bimodality}
  \caption{Distribution of visibility values across individual address
  blocks, comparing the main measurement to the wider one, at their shared
  points in time. Each line is a probability density over 40 equal-width
  bins spanning 0 (fully withdrawn) to 1 (fully visible); the $y$-axis is
  log-scaled so that both the dominant near-0/near-1 modes and the sparse
  middle are visible on one plot (a bar's height, times the bin width of
  0.025, gives the fraction of values falling in that bin -- a height of
  10 on this plot is about 25\% of all values landing in that single bin).
  The shaded band marks the same ``ambiguous middle'' (visibility strictly
  between 0.1 and 0.9) referenced by the \bimodalAmbiguousPrimary{} and
  \bimodalAmbiguousRis{} figures in the text: both measurements place
  almost all of their mass outside it, in an almost identical shape,
  despite drawing on substantially different numbers and locations of
  vantage points.}
  \label{fig:bimodality}
\end{figure}

As an independent check, this paper's routing measurements were compared
against a public routing-history lookup service (RIPEstat) for a sample of
major Iranian networks, and agreed on \ripestatAgreeN{} of \ripestatN{}
network/time-point pairs checked.

\subsection{Measuring reachability and classifying network state}
\label{sec:state-derivation}

Routing visibility alone does not tell us whether a network was actually
reachable -- a network can remain visible in BGP while all traffic to it is
silently blocked. To capture this, this paper also uses IODA, a public
Internet-outage monitoring service that continuously and passively
measures how many hosts in each network respond to background Internet
traffic (an approach that requires no cooperation from, or special access
to, the networks being measured). Each network, at each point in time, is
classified into one of three states: \emph{announced and reachable}
(visible in routing, and responding normally); \emph{announced but dark}
(visible in routing, but the responsiveness signal has collapsed relative
to that network's own normal level -- the signature of traffic filtering
rather than a routing-level shutdown); or \emph{withdrawn} (not visible in
routing at all).

A network's ``normal level'' of responsiveness is its own median value over
a fixed, pre-registered reference month (\probingBaselineStart{} to
\probingBaselineEnd{}, chosen for being distant from both the June 2025
reference blackout and the shutdown itself, before any results were seen).
A network is only included in reachability-based results if it had enough
signal during that reference month to establish a meaningful baseline
(\probingAdequacyMinNonzeroShare{} of the month with a nonzero reading, and
a median of at least \probingAdequacyMinMedian{} units); \nProbingExcluded{}
of \nIrAsns{} networks do not meet this bar and are excluded from
reachability results, though they still contribute to routing-only
results. Because a network with missing or insufficient data defaults to
``reachable'' rather than ``dark'' (classifying a network as dark requires
positive evidence of a responsiveness collapse), every \emph{announced but
dark} figure reported in this paper is a conservative, lower-bound
estimate.

The two numeric thresholds behind this classification -- how visible a
network's routing must be to count as ``announced,'' and how far
responsiveness must fall (relative to a network's own baseline) to count as
``dark'' -- were each set by a dedicated check, run and signed off before any
headline result was treated as final. For the visibility threshold, the
distribution of individual networks' baseline-period visibility values
shows a clean gap between ``clearly visible'' and ``clearly not,'' and the
threshold used (\visibilityAnnouncedMin{}) falls inside that gap
($[\visibilityValleyLo{}, \visibilityValleyHi{}]$); this also happens to
match the definition used by IODA's own methodology (visible to at least
half of monitoring collectors)~\citep{alcock2024ioda}, which was given
priority for comparability with other work once it was confirmed that
either value produces identical results on this paper's data. For the
responsiveness threshold, the value used (\probingDarkRatio{}) is the
largest one for which the control-population networks (which have no
shutdown to detect) are essentially never misclassified as dark, checked
across the entire study period rather than a single reference window so
the result holds up under the full range of ordinary background noise
those networks produce, not just a calm sample of it
(\darkRatioCrossoverRatio{} is the point this stops holding, at a
\darkRatioCrossoverRate{} false-positive rate).

\subsection{Accounting for partial recovery}

During the uneven recovery phase (Restoration, P4), address blocks
sometimes come back online only in part -- a small piece of a larger block
reappearing, rather than the whole thing. To capture this correctly, an
address block's visibility is computed as the fraction of its total
address space actually covered by currently visible announcements, rather
than a simple yes/no check. Where the same address space is announced at
more than one level of specificity at once (which happens routinely in
normal BGP operation), it is only counted once, at the most specific level
announced, so as not to double-count it. A simpler, coarser measure --
whether \emph{any} part of a block has reappeared at all -- is also
reported alongside this primary measure, since it answers a different,
useful question, but it is never used in place of the primary measure
because it cannot tell a fully restored block apart from one with only a
small fragment back online.

\subsection{Comparing shutdown speed across events}

Five historical shutdown or restoration events -- November 2019, June
2025, the January 2026 collapse, the February 2026 onset, and the May 2026
restoration -- are compared at fine (minute-level) time resolution, using
the full stream of
routing announcements and withdrawals during each event window rather than
the coarser periodic snapshots used elsewhere in this paper. Before
comparing events, brief routing flickers (a route withdrawn and
re-announced again within \flapThresholdSeconds{} seconds) are removed, on
the grounds that these are ordinary network noise rather than a real
disconnection. This cutoff was chosen by directly examining the timing
gaps between withdrawals and their following re-announcements across
millions of such events: most reconnect within the first minute in a way
consistent with routine protocol-level retransmission, after which the
pattern changes shape. The exact cutoff is not critical to this paper's
conclusions -- the main comparison across events changes by less than 2\%
whether flickers shorter than 30, 60, or 300 seconds are removed.

The primary metric used to compare how quickly each event unfolded is the
range between the 5th and 95th percentile of when individual address
blocks were first withdrawn, rather than the full first-to-last range. In
practice this means: take the withdrawal time of every affected address
block, sort them, drop the earliest 5\% and the latest 5\%, and measure the
span covered by the remaining middle 90\%. This matters because a handful
of very early or very late outliers -- a network that happened to flap for
an unrelated reason, or one that took days longer than everyone else to be
affected -- can otherwise dominate the full first-to-last range and give a
misleading picture of how the bulk of the event actually unfolded. The May
2026 restoration event makes this concrete: its full first-to-last range
is \durationMayRestorationFullRangeHours{} hours, but the trimmed
5th-95th-percentile range is only \durationMayRestorationRangeHours{}
hours -- a small number of straggling address blocks stretch the untrimmed
range more than fourfold beyond where 90\% of the actual restoration
activity took place. The full, untrimmed range is still reported alongside
the trimmed one throughout, so this trimming is never hidden.

\subsection{Defining ``restored''}

An address block counts as restored at the first point after the
Restoration phase (P4) begins where its visibility (using the
address-space-weighted measure above) climbs back above
\restorationPrimaryThresholdPct{} of that same block's own pre-shutdown
baseline level. Two additional thresholds (25\% and 80\%) are also
reported alongside this one throughout, to show whether the results are
sensitive to exactly where this line is drawn. Separately, how
\emph{completely} a block has recovered by the end of the study is reported
as its average visibility over the final \steadyStateDays{} days, relative
to baseline -- a different question from how \emph{quickly} it recovered,
and reported separately rather than combined into one number. The routine
8-hour measurement snapshots turn out to be too coarse to time recovery
precisely (most blocks appear ``restored'' at the very first snapshot after
the phase boundary, which reflects the measurement's time resolution more
than the true recovery moment), so the fine-grained, minute-level event
data described above is used instead for any result that compares recovery
speed across network categories.

\subsection{How claims are qualified}

Every factual claim in this paper's results and discussion is labeled with
one of three levels of confidence, decided in advance as a matter of
study design rather than chosen sentence-by-sentence after seeing the
results: \emph{observed} (directly shown by the measurement), \emph{strongly
implied} (a pattern in the measurement whose alternative explanations are
implausible, given tight timing and at least one independent, corroborating
source), or \emph{consistent with reporting} (a plausible reading that
depends on outside reports, which the measurement on its own does not
establish). Claims about the intentions or decisions of specific actors --
government orders, ministry decisions, and the like -- are never given
higher confidence than the third level on the basis of this paper's
measurement alone.
